\documentclass{beamer}
% Polish deck (issue #944 leftovers): nested lists, \framesubtitle, sans
% math, \tableofcontents, frame-level and action overlay specifications,
% \temporal, [t]/[b] with allowframebreaks, a footnote inside columns.
\setbeamertemplate{navigation symbols}{}

\title{Polish}
\author{A. Author}
\date{March 2026}

\begin{document}

\begin{frame}
  \frametitle{Outline}
  \tableofcontents
\end{frame}

\section{Lists}

\begin{frame}
  \frametitle{Nested lists}
  \framesubtitle{Three levels of itemize, two of enumerate}
  \begin{itemize}
    \item First level item.
    \begin{itemize}
      \item Second level item, set in small.
      \begin{itemize}
        \item Third level item, set in footnotesize.
      \end{itemize}
      \item Another second level item.
    \end{itemize}
    \item Back to the first level.
  \end{itemize}
  \begin{enumerate}
    \item One
    \begin{enumerate}
      \item One point one
      \item One point two
    \end{enumerate}
    \item Two
  \end{enumerate}
\end{frame}

\section{Mathematics}

\begin{frame}
  \frametitle{Sans-serif mathematics}
  Inline $x + y = 2z$ and $f(x) = \alpha x^2 + \beta$ use the sans family.
  \[
    \int_0^1 f(x)\,dx = \frac{a+b}{2}
  \]
  Operators like $\sin x$ and $\log n$ follow.
\end{frame}

\section{Overlays}

% 2 pages: the frame-level specification suppresses slide 1.
\begin{frame}<2->
  \frametitle{Frame-level specification}
  \begin{itemize}
    \item<1-> Item on every slide.
    \item<2-> Item from the second slide.
    \item<3-> Item from the third slide.
  \end{itemize}
\end{frame}

% 3 pages: an action specification and \temporal.
\begin{frame}
  \frametitle{Actions}
  \begin{itemize}
    \item<1-| alert@2> Alerted on the second slide only.
    \item<2-| alert@3> From the second slide, alerted on the third.
    \item<3-> From the third slide.
  \end{itemize}
  \temporal<2>{Before the second slide.}{On the second slide.}{After the second slide.}
\end{frame}

\section{Breaks}

\begin{frame}[t,allowframebreaks]
  \frametitle{Top-aligned breaks}
  \begin{itemize}
    \item Item one of a list that is too long for a single frame.
    \item Item two.
    \item Item three.
    \item Item four.
    \item Item five.
    \item Item six.
    \item Item seven.
    \item Item eight.
    \item Item nine.
    \item Item ten.
    \item Item eleven.
    \item Item twelve.
    \item Item thirteen.
    \item Item fourteen.
    \item Item fifteen.
    \item Item sixteen.
    \item Item seventeen.
    \item Item eighteen.
    \item Item nineteen.
    \item Item twenty, the last one, on the second page of the frame.
  \end{itemize}
\end{frame}

\begin{frame}[b,allowframebreaks]
  \frametitle{Bottom-aligned breaks}
  \begin{itemize}
    \item Item one of a list that is too long for a single frame.
    \item Item two.
    \item Item three.
    \item Item four.
    \item Item five.
    \item Item six.
    \item Item seven.
    \item Item eight.
    \item Item nine.
    \item Item ten.
    \item Item eleven.
    \item Item twelve.
    \item Item thirteen.
    \item Item fourteen.
    \item Item fifteen.
    \item Item sixteen.
    \item Item seventeen.
    \item Item eighteen.
    \item Item nineteen.
    \item Item twenty, the last one, on the second page of the frame.
  \end{itemize}
\end{frame}

\section{Columns}

\begin{frame}
  \frametitle{A footnote in a column}
  \begin{columns}[T]
    \column{.5\textwidth}
      The left column carries a footnote\footnote{Set at the bottom of the
      frame, not of the column.} in its text.
    \column{.5\textwidth}
      The right column has none.
  \end{columns}
\end{frame}

\end{document}
